\section{Discussion}
\label{sec:discussion}

\paragraph{The verifier needed verifying}
Replacing regular expressions with a parser and range-aware dependency analysis immediately broke the add-on's own test suite: nearly every fixture inserted \code{=SUM(B2:B3)} into B2, a circular reference the shipped verifier never noticed. The real-world corpus then found two defects in our new verifier, and one in our own corpus extraction, in a single pass, and a confirmation set scored with the verifier frozen found a further gap, deleted references, that the benchmark could not have shown because we had not thought of it. A deterministic check is easy to believe; one that is wrong in either direction silently degrades the agent it guards, so a labelled fault benchmark, a clean-negative corpus and real formulas should accompany any such component, and the real formulas should be scored in a way that cannot be tuned after the fact.

\paragraph{What deterministic checks buy, and where they stop}
On the benchmark nearly all loud faults and most silent ones are stopped or flagged, the contribution concentrated in a few cheap layers; on real formulas the two verifiers reject about the same share, and what differs is the kind of fault. End to end the picture is more modest: a few points for an 8B model, about one for a 3B model that can barely act on feedback, and little advantage over the shipped verifier because these models mostly make errors both catch. The hosted models make almost no errors the workbook can see, so the checks are idle and the context decides accuracy; the checks matter for the models that need them, and cost the others nothing. What remains is what the workbook cannot adjudicate (Section~\ref{sec:levels}): a valid formula over the wrong data.

\paragraph{Repair trades visible for invisible failures}
The most useful result may be the least flattering. Verification-guided repair raises accuracy, yet it also raises the number of plausible-but-wrong values written into the sheet, because a formula repaired into validity is no more likely to be right than the one it replaced, and a wrong value that looks right is the failure a user cannot see. In our runs \hmAUnsafeVtwoRate\% (\efAName) to \hmBUnsafeVtwoRate\% (\efBName) of the rejected attempts end that way. Whether that trade is worth it depends on what a visible error costs: a user sees \code{\#REF!} at once, but may never see a wrong total. A design that holds the rejected formula, shows the diagnostics and lets the user decide has no such side effect; one that repairs should follow with an execution check and show the computed value next to the request. The add-on's approval step, dry run and undo address the problem procedurally.

\paragraph{Policy, latent faults and prompts}
Which findings block is a policy: we block only what the workbook proves impossible or a near-miss of something real, and demote plausible-but-probably-wrong findings to \emph{suspicious}, whose repair is worth a further point (Table~\ref{tab:e5-repair}). Value-based checks, including the execution accuracy used by most benchmarks, cannot see about one percent of real formulas (\rwcFGhostOkMasked{} of \rwcFFormulas{} in the confirmation set) that reference a missing sheet behind \code{IFERROR}. And prompts anchor: one of three worked examples in the add-on's prompt uses a sheet called \code{Sheet2}, and on lookup tasks \efAName{} wrote \code{Sheet2!\dots} in \efALkSheetTwo{} of \efALkN{} first attempts although the request names the real sheet. The symbols layer rejects every such formula and its hint lists the real sheets; repair lifts lookup accuracy from \efALkBase\% to \efALkvtwofb\% (resampling: \efALkresample\%). A verifier is no excuse for a careless prompt, but it turns a systematic prompt defect from a silent wrong answer into a one-step repair.

\paragraph{Evaluating a verifier you wrote}
A benchmark written alongside the component it tests measures coverage of the authors' idea of faults. We used four countermeasures, none sufficient alone: labels from an independent engine, real formulas, a confirmation set scored once with a frozen verifier, and a harness for checking the labels in Google Sheets itself (Appendix~\ref{app:gsheets}). The real-formula result was sobering, and it is the one that exposed the gap.
